\section{The Harness}\label{sec:system}

ResearchForge runs an agent (here through DeepSeek Harness, \texttt{dsh}, in headless mode) against a set of typed tools exposed over MCP. The model never writes files directly: every finding is a tool call whose arguments are validated and stored as a record in a per-project workspace. A controller outside the model chooses the next action; each action is one agent turn.

\paragraph{Records and evidence rules.} Records include retrieved papers (merged across arXiv, OpenAlex, Semantic Scholar and Crossref), paper analyses, claims of kind \emph{evidence}, \emph{inference}, \emph{hypothesis} or \emph{assumption}, a landscape, a critique, gaps, uncertainties, directions and experiment plans. A claim may cite only papers retrieved in the same investigation. \emph{Direct} support requires a quote of at least 20 characters that is found verbatim (after Unicode, whitespace and case normalisation; ellipses allowed between fragments) in the stored abstract or the fetched full text; otherwise the tool call is rejected with an explanation. Hypotheses and assumptions never count as grounded.

\paragraph{Controller.} Each iteration the controller (i) raises its own questions, (ii) picks one action, and (iii) measures what changed. A coverage tier runs first: formalize the idea, plan, search until minimum coverage, read, synthesise a landscape, critique. After that, every open uncertainty contributes candidate actions scored by
\[\text{score} = \text{importance}\times\text{expected gain}\times\text{relevance}\times\text{evidence deficiency}\,/\,\text{cost},\]
where expected gain blends a prior per action with the information gain the action actually produced earlier in the run. Information gain counts new relevant papers, analyses, verified claims, resolved uncertainties, contradicted conclusions and new structural records. Investigation closes when the last three steps produced no gain or the best candidate falls below a threshold; hard budgets are safety stops. The final \emph{research decision} is derived by rule from the records, not written by the model: \textsc{already\_well\_explored} if the critique records substantial overlap or the direction is rejected; otherwise \textsc{insufficient\_evidence} if high-importance questions remain unresolved; \textsc{needs\_modification} if the direction changed or the critique judged the idea weak; otherwise \textsc{promising}.

\paragraph{Self-challenge.} When a critique, a gap or a new direction is recorded, the controller raises a high-importance question that asks whether published work contradicts it, to be investigated by a contradiction-seeking search (\textsc{challenge} or \textsc{investigate\_gap}). The agent records a verdict (\emph{holds}, \emph{weakened}, \emph{refuted}, \emph{inconclusive}); \emph{weakened} and \emph{refuted} require contradicting papers that were actually retrieved, and they raise a refinement question. The challenge of the first critique is taken immediately after the critique is recorded, before any other investigation, so that every run with self-challenge performs it at least once regardless of budget.

\paragraph{Critique revision.} If a challenge weakens or refutes a conclusion after the current critique was recorded, the controller re-runs the critique before anything else, so the final assessment never predates evidence against it.

\paragraph{Source integrity.} Bibliographic APIs sometimes attach the wrong abstract to a record; we observed OpenAlex serving another system's abstract for H2O (arXiv 2306.14048). When two sources' abstracts share under 20\% of their distinctive words, the one that fits the title is kept and the record is flagged; an abstract that never names the method in a ``Name: \ldots'' title is flagged. A flagged abstract cannot supply direct evidence; the quote must be found in full text.

\paragraph{Switches.} The components above are independent settings: \texttt{challenge\_conclusions}, \texttt{recritique\_on\_contradiction}, \texttt{integrity\_checks}, and \texttt{controller\_mode}, whose \texttt{pipeline} value runs the coverage tier once, then refine, plan and finalise, with no uncertainty-driven investigation and no self-raised questions. All other behaviour---tools, prompts for shared actions, evidence rules, the decision rule and budgets---is identical across systems.
